\section*{Appendix: Sample data analysis and uncertainty propagation}
\noindent All values below use the reported uncertainty convention: instrument
half-resolution bounds are converted to standard uncertainties by dividing by $\sqrt{3}$,
assuming uniform rounding errors. The uncertainty of locating fringes is assumed to be
$\sigma_N=0.5$ fringe. Common reference errors are not independent observations.

\subsection*{Green laser: wavelength from fitted slope}
\begin{align*}
N &= m_\lambda\Delta x+b,\\
\lambda &=\frac{2000}{m_\lambda}\;\mathrm{nm},\\
\sigma_\lambda &=\left|\frac{\partial\lambda}{\partial m_\lambda}\right|\sigma_m
=\frac{2000\sigma_m}{m_\lambda^2}.
\end{align*}
For green trial 1, $m_\lambda=3.60253\pm0.07300$ fringes/$\mu$m, hence
$\lambda_1=555.17\pm11.25$ nm.

\subsection*{Measurement uncertainty in mirror displacement}
If each micrometer reading has a maximum rounding error of $0.5\,\mu$m,
$\sigma_{x_i}=0.5/\sqrt{3}\,\mu$m. For a difference
$\Delta x_i=x_i-x_0$, the standard uncertainty of a single difference is
\[
\sigma_{\Delta x}=\sqrt{\sigma_{x_i}^2+\sigma_{x_0}^2}
=\sqrt{2}\frac{0.5}{\sqrt3}\,\mu\mathrm m.
\]
The same $x_0$ is shared by the data, so the point-to-point covariance is not zero.
A free intercept absorbs a constant reference offset when estimating the slope.

\subsection*{Angular transformation and refractive index}
\begin{align*}
\theta &= \frac{\pi}{180}\theta_{\mathrm{deg}},\qquad u=\theta^2,\\
\sigma_u&=\left|\frac{du}{d\theta}\right|\sigma_\theta=2|\theta|\sigma_\theta,\\
n&=\frac{1}{1-m_n\lambda/t}.
\end{align*}
At $\theta=1^\circ$, $u=0.0003046$ rad$^2$ and with a single-angle
standard uncertainty of $0.0050383$ rad, $\sigma_u=0.0001759$ rad$^2$.

For independent slope, wavelength, and thickness uncertainties,
\[
\sigma_n=n^2\sqrt{\left(\frac{\lambda\sigma_m}{t}\right)^2+
\left(\frac{m\sigma_\lambda}{t}\right)^2+
\left(\frac{m\lambda\sigma_t}{t^2}\right)^2}.
\]
For the 134$^\circ$+ direction $m_n=6678.853\pm272.118$ rad$^{-2}$,
$t=7790000$ nm, and $\lambda=550.710$ nm. This gives $n=1.8945$ and
$\sigma_n=0.0718$. The shared-reference-angle
bias is an additional possible systematic error, treated separately.

\subsection*{Temperature difference and aluminium expansion}
\begin{align*}
\Delta T_i &= T_i-T_0,\\
\sigma_{\Delta T_i}&=\sqrt{\sigma_{T_i}^2+\sigma_{T_0}^2},\\
\alpha &= \frac{m_\alpha\lambda}{2L_0},\\
\frac{\sigma_\alpha}{|\alpha|} &=
\sqrt{\left(\frac{\sigma_{m_\alpha}}{m_\alpha}\right)^2+
\left(\frac{\sigma_\lambda}{\lambda}\right)^2+
\left(\frac{\sigma_{L_0}}{L_0}\right)^2}.
\end{align*}
At $T=27.0^\circ$C, $T_0=25.2^\circ$C and each reading has
$\sigma_T=0.05/\sqrt3$ C, the difference is $\Delta T=1.8^\circ$C with
$\sigma_{\Delta T}=0.0408^\circ$C.
The fitted slope is $m_\alpha=6.95191\pm0.04032$ fringes/$^\circ$C,
giving $\alpha=(21.479\pm0.280)\times10^{-6}\,^\circ\mathrm{C}^{-1}$.

\subsection*{Goodness of fit}
\[
\chi^2=\sum_i\frac{(y_i-\widehat y_i)^2}{\sigma_{y_i}^2},
\qquad \chi^2_\nu=\frac{\chi^2}{N-p}
\]
The above conventional expression applies to independent errors in $y$ only.
Our primary two-axis fit minimizes squared normalized corrections in both
$x$ and $y$, and reports the corresponding total $\chi^2$ divided by $N-p$.
$R^2=1-\mathrm{SSE}/\mathrm{SST}$ is computed from vertical residuals at
recorded x-values. Exact reference-angle and thermal-lag systematics are not
captured by these reported statistical uncertainties.
\subsection*{Course handout: OLS slope, uncertainty, and residual}
\[
m=\frac{N\sum_i x_i y_i-(\sum_i x_i)(\sum_i y_i)}{N\sum_i x_i^2-(\sum_i x_i)^2},\qquad
s_m=\sqrt{\frac{Ns_{y,x}^2}{\Delta}},\qquad
s_{y,x}^2=\frac{\sum_i(y_i-mx_i-b)^2}{N-2}.
\]
For green trial 1, $N=5$ and $\Delta=1250$. The ordinary least-squares slope is $m=3.6000\pm0.0611$ fringes/$\mu$m.
At $\Delta x=5.0\,\mu$m, the observed fringe count is $15$, the prediction is $14.200$, and residual is $0.800$ fringes.
The effective vertical residual-bar uncertainty in a linear fit is $\sigma_{r,\mathrm{eff}}\simeq\sqrt{\sigma_N^2+(m\sigma_{x,\mathrm{ind}})^2}$.

